\scanpage{33}
\mspage{16}
6. \uline{Relation} \uline{avec l'octaèdre non orienté} $\Delta_M$.

Soit

\[
\begin{aligned}
M^*&=\{x\in M\mid x^0\,(\text{image dans }M_0)\ne0\}\\
&=\text{image inverse de }J=M_0^*\subset M_0\text{ par } M\longrightarrow M_0 .
\end{aligned}\tag*{(29)}
\]

\deleted{\unsure{donc on|donc on a}} donc $M^*$ est un torseur relatif sur $M_0^{\unsure{a|*|\#}}=J$, du groupe $(M_0)_J$ (i.e famille constante de groupes isom. à $M_0$, indexée par $J$). Le groupe $(\mathbb Z/4\mathbb Z)^*=\{\pm1\}$ opère sur $M^{*}_{\deleted{\unsure{x|0}}}$, de façon compatible avec la projection $M^*\deleted{\unsure{/\,\{\}|\to}}\to J$, on pose

\[
\Delta_M=M^*/(\mathbb Z/4\mathbb Z)^*=M^*/\{\pm1\}=\mathbb P(M)(\mathbb Z/4\mathbb Z)\tag*{(30)}
\]

ce sont les points \deleted{\unsure{sur}} à val. dans $\Lambda=\mathbb Z/4\mathbb Z$ de la droite projective sur $\Lambda$ définie par $M$, et l'application \deleted{\unsure{canonique}} \inserted{can.}

\[
\Delta_M\longrightarrow J \tag*{(31)}
\]

s'identifie à l'application de restriction, ou réduction mod $2$

\[
\mathbb P(M)(\mathbb Z/4\mathbb Z)\longrightarrow\mathbb P(M)(\mathbb F_2).
\]

C'est une application de degré $2$, qui définit donc sur $\Delta$ la structure de (l'ens. des sommets) d'un octaèdre, \deleted{dont} $J$ est l'ens des diagonales. Les \deleted{Le groupe $\mathrm{Gl}(M)$ opère sur} \inserted{12 arêtes de l'octaèdre sont les} éléments de $\mathcal P_2(\Delta)$ (qui a $\frac{6\cdot5}{2}=15$ éléments) qui ne sont pas \deleted{\unsure{des diago|des diagon|des diagonales}} l'une des trois diagonales, i.e. des trois fibres de (31). Les $8$ faces de l'octaèdre sont les parties à trois éléments $\inserted{f}\deleted{\unsure{Q|D|6}}$ de $\Delta$ telles que l'application induite par (31) soit une bijection

\[
f\xrightarrow{\sim}J\ ; \tag*{(32)}
\]

Elles s'identifient donc aux sections de (31), tandis que les arêtes s'identifient aux sections partielles sur des parties de $J$ à deux éléments, et bien
% archive sheet number: 32
